\slajdid{07-s018}
\begin{frame}[label=07-s018]
\gusttekst\centering\begin{tikzpicture}[font=\fontsize{8}{9}\selectfont,>=Latex]
\node[draw,fill=DEplava!10,minimum height=.45cm] (entry) at (0,0) {Design Entry};
\node[draw,fill=DEplava!10,minimum height=.45cm] (test) at (4.5,0) {Test Development};
\node[draw,fill=DEplava!10,minimum height=.45cm] (syn) at (0,-.8) {Synthesis};
\node[draw,fill=DEplava!10,minimum height=.45cm] (func) at (3.2,-.8) {Functional Simulation};
\node[draw,fill=DEplava!10,minimum height=.45cm] (map) at (0,-1.6) {Device Mapping};
\node[draw,fill=DEplava!10,minimum height=.45cm] (time) at (4.5,-1.6) {Timing Simulation};
\draw[->,line width=1.2pt,DEplava] (entry)--(syn);\draw[->,line width=1.2pt,DEplava] (entry.east)-|([xshift=-4mm]func.north);
\draw[->,line width=1.2pt,DEplava] (test.south)--++(0,-.12)-|([xshift=4mm]func.north);\draw[->,line width=1.2pt,DEplava] (test.east)--(6,0)|-(time.east);
\draw[->,line width=1.2pt,DEplava] (syn)--(map);\draw[->,line width=1.2pt,DEplava] (map)--(time);
\draw[->,line width=1.2pt,DEplava] (map)--(0,-2.18);
\filldraw[fill=orange!35,draw=DEplava] (-.85,-2.6)--(.6,-2.6)--(.95,-2.22)--(-.5,-2.22)--cycle;
\draw[DEplava] (-.85,-2.6)--(-.85,-2.66)--(.6,-2.66)--(.95,-2.28)--(.95,-2.22);
\draw[DEplava] (-0.85,-2.66)--(-0.85,-2.84);
\draw[DEplava] (-0.705,-2.66)--(-0.705,-2.84);
\draw[DEplava] (-0.56,-2.66)--(-0.56,-2.84);
\draw[DEplava] (-0.41500000000000004,-2.66)--(-0.41500000000000004,-2.84);
\draw[DEplava] (-0.27,-2.66)--(-0.27,-2.84);
\draw[DEplava] (-0.125,-2.66)--(-0.125,-2.84);
\draw[DEplava] (0.019999999999999907,-2.66)--(0.019999999999999907,-2.84);
\draw[DEplava] (0.16499999999999992,-2.66)--(0.16499999999999992,-2.84);
\draw[DEplava] (0.30999999999999994,-2.66)--(0.30999999999999994,-2.84);
\draw[DEplava] (0.45499999999999996,-2.66)--(0.45499999999999996,-2.84);
\draw[DEplava] (0.6,-2.66)--(0.6,-2.84);
\draw[DEplava] (.9,-2.33)--(.9,-2.51);
\end{tikzpicture}
\par\vspace{2mm}\raggedright
Gde se VHDL uklapa u ovaj dijagram? VHDL (zajedno sa drugim oblicima unosa, kao što su šeme i blok dijagrami) će se koristiti za unos dizajna. VHDL izvorni kod se može uneti u simulaciju, omogućavajući mu da bude funkcionalno verifikovan, ili može biti prosleđen direktno alatima za sintezu za implementaciju u određenom tipu uređaja. Na strani razvoja testa, mogu se kreirati VHDL testbench koje testiraju kolo kako bi se proverilo da li ispunjava funkcionalna i vremenska ograničenja specifikacije. Testbench se mogu uneti pomoću editora teksta ili mogu biti generisani iz drugih oblika informacija o test stimulansima, kao što su grafički talasni oblici. Za preciznu vremensku simulaciju kola posle rutiranja, verovatno ćete koristiti program za generisanje vremenskog modela dobijen od dobavljača uređaja ili treće strane dobavljača simulacionog modela. Alati za generisanje modela, kao što je ovaj, obično generišu VHDL izvorne datoteke sa vremenskim napomenama koje podržavaju veoma precizne simulacije na nivou sistema.
\end{frame}
